% RQ5 Background: Inter-DAO Cooperation

\subsection{Inter-DAO Cooperation in Practice}

\subsubsection{Protocol Partnerships}

DAOs increasingly form partnerships:
\begin{itemize}
    \item \textbf{Curve-Convex}: Deep integration of governance and incentives
    \item \textbf{Lido-Aave}: Coordinated staking and lending mechanisms
    \item \textbf{Optimism Collective}: Bicameral structure linking Token House and Citizens' House
\end{itemize}

\subsubsection{Grants Coordination}

Multiple DAOs co-fund public goods:
\begin{itemize}
    \item \textbf{Protocol Guild}: Multi-protocol funding for Ethereum core developers
    \item \textbf{Gitcoin Grants}: Quadratic funding with matching from multiple treasuries
    \item \textbf{Ecosystem funds}: Joint grant programs across aligned protocols
\end{itemize}

\subsubsection{Standard Setting}

Cross-DAO governance of shared standards:
\begin{itemize}
    \item \textbf{ERC standards}: Informal multi-community consensus
    \item \textbf{Metagovernance}: DAOs holding and voting tokens of other DAOs
\end{itemize}

\subsection{Challenges of Cooperation}

\subsubsection{Preference Divergence}

DAOs may have conflicting preferences:
\begin{itemize}
    \item Competitive dynamics (market share, user acquisition)
    \item Different time horizons (short-term vs. long-term focus)
    \item Ideological differences (decentralization philosophy, value capture)
\end{itemize}

\subsubsection{Coordination Costs}

Cooperation imposes overhead:
\begin{itemize}
    \item Synchronization of governance timelines
    \item Cross-community communication
    \item Dispute resolution mechanisms
\end{itemize}

\subsubsection{Free Riding}

Public goods created through cooperation benefit all, creating free-rider incentives:
\begin{itemize}
    \item Under-contribution to joint initiatives
    \item Reliance on others to fund common goods
    \item Exit when costs exceed perceived benefits
\end{itemize}

\subsection{Theoretical Foundations}

\subsubsection{Collective Action Theory}

Inter-DAO cooperation faces classic collective action problems \citep{olson1965logic}:
\begin{itemize}
    \item Large groups face coordination challenges
    \item Individual incentives may diverge from collective optimum
    \item Selective incentives or small groups enable cooperation
\end{itemize}

\subsubsection{Commons Governance}

\citet{ostrom1990governing} identified conditions for successful commons management:
\begin{itemize}
    \item Clearly defined boundaries
    \item Congruent rules and local conditions
    \item Collective-choice arrangements
    \item Monitoring and graduated sanctions
\end{itemize}

These principles apply to inter-DAO coordination of shared resources.

\subsubsection{Coalition Formation}

Game-theoretic coalition models inform multi-DAO dynamics:
\begin{itemize}
    \item Coalition stability (no subset wants to deviate)
    \item Fair division of coalition benefits
    \item Blocking coalitions and veto power
\end{itemize}

\subsection{Prior Work}

Limited academic work addresses inter-DAO cooperation specifically. \citet{tan2023agentdao} explored LLM-based agent DAOs with emergent coordination. DeepDAO \citep{deepdao2023analytics} tracks cross-DAO token holdings. Our simulation framework provides systematic analysis of cooperation dynamics.
